\chapter{TOROIDAL AND POLOIDAL VECTOR FIELDS}
\label{app:III}

\section{A general characterization of solenoidal vector fields. The fundamental basis}\label{sec:128}

\textsc{The} vector fields defined by the equations
\begin{equation}
\label{eq:A3-1}
\mathbf{T} = \operatorname{curl}\!\left(\frac{\Psi}{r}\,\mathbf{r}\right)
\end{equation}
and
\begin{equation}
\label{eq:A3-2}
\mathbf{S} = \operatorname{curl}\!\left[\operatorname{curl}\!\left(\frac{\Phi}{r}\,\mathbf{r}\right)\right],
\end{equation}
where $\Psi$ and $\Phi$ are arbitrary scalar functions of position, are clearly solenoidal. A field derivable in the manner~(\ref{eq:A3-1}) is said to be \emph{toroidal} with the \emph{defining scalar} $\Psi$; and a field derivable in the manner~(\ref{eq:A3-2}) is said to be \emph{poloidal} with the defining scalar $\Phi$. Alternative ways of writing $\mathbf{T}$ and $\mathbf{S}$ are
\begin{equation}
\label{eq:A3-3}
\mathbf{T} = \operatorname{grad}\frac{\Psi}{r} \times \mathbf{r}
\end{equation}
and
\begin{equation}
\label{eq:A3-4}
\mathbf{S} = \operatorname{curl}\!\left(\operatorname{grad}\frac{\Phi}{r} \times \mathbf{r}\right).
\end{equation}

In a system of spherical coordinates, $(r, \vartheta, \varphi)$, the components of $\mathbf{S}$ and $\mathbf{T}$ are:
\begin{equation}
\label{eq:A3-5}
T_r = 0, \quad T_\vartheta = \frac{1}{r\sin\vartheta}\frac{\partial\Psi}{\partial\varphi}, \quad T_\varphi = -\frac{1}{r}\frac{\partial\Psi}{\partial\vartheta};
\end{equation}
and
\begin{equation}
\label{eq:A3-6}
S_r = \frac{1}{r^2}\,L^2\Phi, \quad S_\vartheta = \frac{1}{r}\frac{\partial^2\Phi}{\partial r\partial\vartheta}, \quad S_\varphi = \frac{1}{r\sin\vartheta}\frac{\partial^2\Phi}{\partial r\partial\varphi},
\end{equation}
where $L^2$ stands for the operator (defined in equation~VI~(25))
\begin{equation}
\label{eq:A3-7}
L^2 = -\frac{1}{\sin\vartheta}\frac{\partial}{\partial\vartheta}\sin\vartheta\frac{\partial}{\partial\vartheta} - \frac{1}{\sin^2\vartheta}\frac{\partial^2}{\partial\varphi^2}.
\end{equation}

Since $\mathbf{T}$ is solenoidal,
\begin{equation}
\label{eq:A3-8}
(\operatorname{curl}^2\mathbf{T})_i = -\frac{\partial^2 T_i}{\partial x_s^2} = \frac{\partial^2}{\partial x_s^2}\,\epsilon_{ijk}\,x_j\frac{\partial}{\partial x_k}\frac{\Psi}{r} = \epsilon_{ijk}\,x_j\frac{\partial}{\partial x_k}\nabla^2\frac{\Psi}{r};
\end{equation}
thus,
\begin{equation}
\label{eq:A3-9}
\operatorname{curl}^2\mathbf{T} = -\operatorname{grad}\!\left(\nabla^2\frac{\Psi}{r}\right) \times \mathbf{r}.
\end{equation}

Similarly,
\begin{equation}
\label{eq:A3-10}
\operatorname{curl}\,\mathbf{S} = -\operatorname{grad}\!\left(\nabla^2\frac{\Phi}{r}\right) \times \mathbf{r}.
\end{equation}

From equations~(\ref{eq:A3-3}), (\ref{eq:A3-4}), (\ref{eq:A3-9}), and~(\ref{eq:A3-10}) it is apparent that $\operatorname{curl}\,\mathbf{T}$ is a poloidal field with the same defining scalar $\Psi$, while $\operatorname{curl}^2\mathbf{T}$ is a toroidal field with the defining scalar
\begin{equation}
\label{eq:A3-11}
\hat{\Psi} = -r\nabla^2\frac{\Psi}{r};
\end{equation}
correspondingly, $\operatorname{curl}\,\mathbf{S}$ is a toroidal field with the defining scalar
\begin{equation}
\label{eq:A3-12}
\hat{\Phi} = -r\nabla^2\frac{\Phi}{r},
\end{equation}
while $\operatorname{curl}^2\mathbf{S}$ is a poloidal field with the same defining scalar $\hat{\Phi}$.

We obtain a \emph{fundamental} basis (on a sphere) for these toroidal and poloidal fields by expanding $\Psi$ and $\Phi$ in spherical harmonics with coefficients which are allowed to be functions of $r$. Thus, the elements of the basis are the fields derived from the scalars
\begin{equation}
\label{eq:A3-13}
\Psi = T(r) Y_l^m(\vartheta, \varphi) \quad \text{and} \quad \Phi = S(r) Y_l^m(\vartheta, \varphi),
\end{equation}
where
\begin{equation}
\label{eq:A3-14}
Y_l^m(\vartheta, \varphi) = e^{im\varphi} P_l^m(\cos\vartheta).
\end{equation}

The components of the fields defined by these special scalars are:
\begin{equation}
\label{eq:A3-15}
T_r = 0, \quad T_\vartheta = \frac{T(r)}{r\sin\vartheta}\frac{\partial Y_l^m}{\partial\varphi}, \quad T_\varphi = -\frac{T(r)}{r}\frac{\partial Y_l^m}{\partial\vartheta},
\end{equation}
and
\begin{equation}
\label{eq:A3-16}
S_r = \frac{l(l+1) S(r)}{r^2}\,Y_l^m, \quad S_\vartheta = \frac{1}{r}\frac{dS}{dr}\frac{\partial Y_l^m}{\partial\vartheta}, \quad S_\varphi = \frac{1}{r\sin\vartheta}\frac{dS}{dr}\frac{\partial Y_l^m}{\partial\varphi},
\end{equation}
where in simplifying the expression for $S_r$, we have made use of the identity (equation~VI~(28))
\begin{equation}
\label{eq:A3-17}
L^2 Y_l^m = l(l+1) Y_l^m.
\end{equation}

The defining scalars of the associated fields $\operatorname{curl}^n\mathbf{T}$ and $\operatorname{curl}\,\mathbf{S}$ are readily found. We have (cf.\ equation~VI~(32))
\begin{equation}
\label{eq:A3-18}
\hat{\Psi} = -r Y_l^m \left(\frac{d^2}{dr^2} + \frac{2}{r}\frac{d}{dr} - \frac{l(l+1)}{r^2}\right)T = -r Y_l^m \left[\frac{1}{r}\frac{d^2}{dr^2}(rT) - \frac{l(l+1)}{r^2}T\right]
\end{equation}
and similarly,
\begin{equation}
\label{eq:A3-19}
\hat{\Phi} = -r Y_l^m \left[\frac{1}{r}\frac{d^2}{dr^2}(rS) - \frac{l(l+1)}{r^2}S\right].
\end{equation}

\section{The orthogonality properties of the basic toroidal and poloidal fields}\label{sec:129}

It is convenient to define on the unit sphere the vector operator
\begin{equation}
\label{eq:A3-20}
\mathbf{V}_r = \mathbf{1}_\vartheta \frac{\partial}{\partial\vartheta} + \frac{1}{\sin\vartheta}\,\mathbf{1}_\varphi \frac{\partial}{\partial\varphi},
\end{equation}
where $\mathbf{1}_\vartheta$ and $\mathbf{1}_\varphi$ are unit vectors, along the principal meridian and the latitude circle, through the point considered. The effect of $\mathbf{V}_r$ on any scalar function $\Psi(r, \vartheta, \varphi)$ is to yield the vector
\begin{equation}
\label{eq:A3-21}
\mathbf{V}_r \Psi = \mathbf{1}_\vartheta \frac{\partial\Psi}{\partial\vartheta} + \frac{1}{\sin\vartheta}\,\mathbf{1}_\varphi \frac{\partial\Psi}{\partial\varphi}.
\end{equation}

On the other hand, if
\begin{equation}
\label{eq:A3-22}
\boldsymbol{\xi} = \mathbf{1}_\vartheta \xi_\vartheta + \mathbf{1}_\varphi \xi_\varphi
\end{equation}
is a two-dimensional vector point-function (on the unit sphere), then we define
\begin{equation}
\label{eq:A3-23}
\operatorname{div}\boldsymbol{\xi} = \mathbf{V}_r \cdot \boldsymbol{\xi} = \frac{1}{\sin\vartheta}\frac{\partial}{\partial\vartheta}(\sin\vartheta\,\xi_\vartheta) + \frac{1}{\sin\vartheta}\frac{\partial\xi_\varphi}{\partial\varphi}.
\end{equation}

Accordingly,
\begin{equation}
\label{eq:A3-24}
\nabla^2_\vartheta\Psi = \mathbf{V}_r \cdot (\mathbf{V}_r\Psi) = \mathbf{V}_r \cdot \left( \mathbf{1}_\vartheta \frac{\partial\Psi}{\partial\vartheta} + \frac{1}{\sin\vartheta}\,\mathbf{1}_\varphi \frac{\partial\Psi}{\partial\varphi} \right) = -L^2\Psi.
\end{equation}

Hence, in particular,
\begin{equation}
\label{eq:A3-25}
\nabla^2_\vartheta Y_l^m = -L^2 Y_l^m = -l(l+1) Y_l^m.
\end{equation}

Two elementary consequences of the foregoing definitions may be noted. \emph{First}, if $\boldsymbol{\xi}$ is a single-valued vector point-function on the unit sphere, then
\begin{equation}
\label{eq:A3-26}
\int\!\!\int \mathbf{V}_r \cdot \boldsymbol{\xi} \, d\Sigma = \int\!\!\int \left[ \frac{1}{\sin\vartheta}\frac{\partial}{\partial\vartheta}(\xi_\vartheta \sin\vartheta) + \frac{1}{\sin\vartheta}\frac{\partial\xi_\varphi}{\partial\varphi} \right] \sin\vartheta \, d\vartheta \, d\varphi
\end{equation}
\[
= \int_0^\pi d\vartheta (\xi_\vartheta \sin\vartheta) + \int_0^{2\pi} d\varphi (\xi_\varphi) = 0.
\]

Thus, in the present geometry Gauss's theorem takes the form
\begin{equation}
\label{eq:A3-27}
\int\!\!\int \mathbf{V}_r \cdot \boldsymbol{\xi} \, d\Sigma = 0.
\end{equation}

A second consequence which directly follows from the definitions is
\begin{equation}
\label{eq:A3-28}
\mathbf{V}_r \cdot (\Psi\boldsymbol{\xi}) = \Psi(\mathbf{V}_r \cdot \boldsymbol{\xi}) + \boldsymbol{\xi} \cdot (\mathbf{V}_r\Psi).
\end{equation}

Making use of the results expressed in equations~(\ref{eq:A3-27}) and~(\ref{eq:A3-28}), we deduce the following important identity:
\begin{equation}
\label{eq:A3-29}
0 = \int\!\!\int \mathbf{V}_r \cdot (\Psi'\mathbf{V}_r\Psi'') \, d\Sigma
\end{equation}
\[
= \int\!\!\int (\Psi'\nabla^2_\vartheta\Psi'' + (\mathbf{V}_r\Psi') \cdot (\mathbf{V}_r\Psi'')) \, d\Sigma
\]
\[
= -l(l+1) \int\!\!\int \Psi' \Psi'' \, d\Sigma + \int\!\!\int (\mathbf{V}_r\Psi') \cdot (\mathbf{V}_r\Psi'') \, d\Sigma.
\]

From the known orthogonality properties of the spherical harmonics, we can infer that
\begin{equation}
\label{eq:A3-30}
\int\!\!\int (\mathbf{V}_r Y_l^m) \cdot (\mathbf{V}_r Y_{l'}^{m'}) \, d\Sigma = l(l+1) N_l^m \delta_{ll'} \delta_{mm'},
\end{equation}
where
\begin{equation}
\label{eq:A3-31}
N_l^m = \frac{4\pi}{2l+1}\frac{(l+|m|)!}{(l-|m|)!}.
\end{equation}

The various orthogonality relations among the basic toroidal and poloidal fields enumerated in \S~46 (p.~226) are now seen to be direct consequences of equation~(\ref{eq:A3-30}). Thus, considering two toroidal fields, $\mathbf{T}_l^m$ and $\mathbf{T}_{l'}^{m'}$, associated, respectively, with the spherical harmonics $Y_l^m$ and $Y_{l'}^{m'}$, we have
\begin{equation}
\label{eq:A3-32}
r^2 \int\!\!\int \mathbf{T}_l^m \cdot \mathbf{T}_{l'}^{m'} \, d\Sigma = T_l T_{l'} \int\!\!\int \left\{ \frac{1}{\sin^2\vartheta}\frac{\partial Y_l^m}{\partial\varphi}\frac{\partial Y_{l'}^{m'}}{\partial\varphi} + \frac{\partial Y_l^m}{\partial\vartheta}\frac{\partial Y_{l'}^{m'}}{\partial\vartheta} \right\} d\Sigma
\end{equation}
\[
= T_l T_{l'} \int\!\!\int (\mathbf{V}_r Y_l^m) \cdot (\mathbf{V}_r Y_{l'}^{m'}) \, d\Sigma
\]
\[
= l(l+1) N_l^m T_l T_{l'} \, \delta_{ll'} \delta_{mm'};
\]
and similarly,
\begin{equation}
\label{eq:A3-33}
r^2 \int\!\!\int \mathbf{S}_l^m \cdot \mathbf{S}_{l'}^{m'} \, d\Sigma = l(l+1)\,l'(l'+1)\frac{S_l S_{l'}}{r^2} \int\!\!\int Y_l^m Y_{l'}^{m'} \, d\Sigma +
\end{equation}
\[
+ \frac{dS_l}{dr}\frac{dS_{l'}}{dr}\int\!\!\int (\mathbf{V}_r Y_l^m) \cdot (\mathbf{V}_r Y_{l'}^{m'}) \, d\Sigma
\]
\[
= l(l+1) N_l^m \left[ \frac{l(l+1)}{r^2}\,S_l S_{l'} + \frac{dS_l}{dr}\frac{dS_{l'}}{dr} \right] \delta_{ll'}\delta_{mm'}.
\]

And finally, the orthogonality of every poloidal field to every toroidal field can be established as follows:
\begin{equation}
\label{eq:A3-34}
r^2 \int\!\!\int \mathbf{S}_l^m \cdot \mathbf{T}_{l'}^{m'} \, d\Sigma = \frac{T_{l'}}{r}\frac{dS_l}{dr} \int_0^{2\pi}\!\!\int_0^\pi \left\{ \frac{\partial Y_l^m}{\partial\vartheta}\frac{\partial Y_{l'}^{m'}}{\partial\varphi} - \frac{\partial Y_l^m}{\partial\varphi}\frac{\partial Y_{l'}^{m'}}{\partial\vartheta} \right\} d\vartheta\, d\varphi
\end{equation}
\[
= \frac{T_{l'}}{r}\frac{dS_l}{dr} \cdot 2\pi i \sin\vartheta_m \left[P_l^{|m|}(\vartheta)\,P_{l'}^{|m'|}(\vartheta)\right]_{\vartheta=0}^{\vartheta=\pi}.
\]

The quantity in square brackets in the last line of~(\ref{eq:A3-34}) clearly vanishes if $l$ and $l'$ are both even or both odd; but if only one of them, say $l$, is odd, then, $P_l^m(\pm 1)$ vanishes at $\theta = 0$ and $\theta = \pi$; hence, in all cases the quantity in square brackets vanishes; and the required orthogonality property is established.

\section*{BIBLIOGRAPHICAL NOTES}

To the references given in the Bibliographical Notes for Chapter~VI (Nos.~13 and~14, p.~269) we may add the following:
\begin{enumerate}
\item J.~M. \textsc{Blatt} and V.~F. \textsc{Weisskopf}, \textit{Theoretical Nuclear Physics}, Appendix~B, John Wiley \& Sons, New York, 1952.
\item G.~E. \textsc{Backus}, `A class of self-sustaining dissipative spherical dynamos', \textit{Annals of Physics}, \textbf{4}, 372--447 (1958).
\end{enumerate}
